\chapter{«عشرون ثانية»}
\label{sec:1}
قال ضابط استخبارات لصحفي، وهو يصف دوره في خط إنتاج للأهداف مبنيّ حول قوائم قتل تولّدها الآلة، إنه كان يقضي عشرين ثانية على كل اسم، وهي مدة تكفي للتحقق من أن الهدف رجل: «لا يستغرق الأمر وقتًا طويلًا لتعرف إن كان صوت الشخص صوت رجل أم امرأة». وأضاف أنه لم تكن له أي قيمة مضافة بوصفه إنسانًا، سوى أنه ختم موافقة. «لقد وفّر ذلك الكثير من الوقت» \autocite{abraham2024lavender}. وفي روايته، هذا ما آل إليه «الإنسان في الحلقة» تحت إيقاع العمليات.\footnote{تتنبأ أبحاث انحياز الأتمتة بهذا التفويض: حين تتنافس مطالب أخرى على انتباههم، يتحقق الأشخاص الذين تُقدَّم لهم توصية آلة منها أقل ويُذعنون لها أكثر \autocite{parasuraman2010complacency}.}

أحد التفسيرات أن الضابط ذاتٌ لم يُمنح ما يكفي من الوقت، وأن العلاج هو أن نُحسن حكمه — مزيد من الوقت، ومراجع ثانٍ، وضمير إضافي يُركَّب في الحلقة. فمن منظور \emph{الإخضاع الاجتماعي} هذا، الإنسان ذاتٌ تقف خارج الآلة وتشغّلها بوصفها أداة: سائق وسيارة، عامل ومخرطة. للإنسان دور، واسم، ومسؤوليات، وضمير يمكن مخاطبته. أما المنظور البديل فهو \emph{الاستعباد الآلي}، حيث الإنسان مكوّن من مكوّنات الآلة، يُجمَع مع أجزائها الأخرى في نظام يضبطها معًا؛ تستخدم الآلة الإنسان عضوًا من أعضائها، وما يسهم به الإنسان وظيفةٌ يؤديها بمعدل يحدده النظام.\footnote{يسمّي دوكتورو موقع الضابط «القنطور المعكوس»: ليس إنسانًا تساعده آلة، بل مكوّنًا لآلة مصنوعًا من إنسان \autocite{doctorow2026reversecentaur}.} والمنظوران لا يستبعد أحدهما الآخر؛ فقد يشغل المرء دورًا في مستوى بينما يكون جزءًا في المستوى الآخر \autocite[الهضبة 13، «The Apparatus of Capture»: يكون هناك استعباد حين يكون البشر أنفسهم قطعًا مكوِّنة لآلة يؤلّفونها فيما بينهم ومع أشياء أخرى، تحت سيطرة وحدة عليا وتوجيهها، ويكون هناك إخضاع حين تُنشئ الوحدةُ العليا الإنسانَ ذاتًا مرتبطة بموضوع خارجي جديد، قد يكون حيوانًا أو أداة أو حتى آلة]{deleuze1987plateaus}، والضابط صاحب الثواني العشرين كان الاثنين معًا. بوصفه جزءًا، كان الموضع الذي يُنتج فيه خط الإنتاج ختمه، وكانت الثواني العشرون زمن دورة ذلك الجزء. وبوصفه ذاتًا، كان مخاطَبًا طوال الوقت: فقد وصف لصحفي التنافر بين أخبار مقتل المدنيين والضغط في العمل لتقديم مزيد من الأهداف. وحتى حين بلغ التنافر ضميره، مضى مع التيار. لقد أحسّ به، ولم يحسم شيئًا؛\footnote{أو على الأقل، لا شيء منع جيش الدفاع الإسرائيلي من قتل 100,000 شخص في غزة بين أكتوبر 2023 وأكتوبر 2025 \autocite{gomezugarte2025uncertainty,acosta2026gaza,economist2025gaza}. وقد نُشرت رواية ذلك الضابط بالتحديد في أبريل 2024 في +972 Magazine. ولا أحد يعرف إن كان قرار ذلك الضابط التحدث إلى صحفي قد منع وفاة إضافية واحدة، أو 50,000 وفاة إضافية، أو لم يمنع أيًّا منها. وماذا كان سيحدث لو رفض الضابط بعد أسبوعين في منصبه — لنقل، افتراضًا، في نوفمبر 2023؟ نحن نتحدث عن ضابط استخبارات يستخدم Lavender لإعداد قوائم القتل. وكانت بيئة عمله اليومية على الأرجح أقرب إلى وظيفة مكتبية مدنية منها إلى فيلم حركة. هل كان عليه أن يطلب اجتماعات متكررة مع رؤسائه، صعودًا في التسلسل القيادي، وهو يرفض المصادقة على أهداف لغارات جوية كان من المعقول توقّع أنها ستسبب وفيات بين المدنيين؟ يمكن القول إن ذلك كان التصرف الأخلاقي. وفي هذا الافتراض المضاد، هل كان القتلى سيكونون أقل؟} أما الحكم بأن السياسة كانت غير متناسبة فجاء، كما قال، بأثر رجعي.

إذن، ما كان ينقص لم يكن الضمير؛ فقد كان لديه ضمير. يعمل الضمير بأن يجعل للقرار ثمنًا، لكن الأثمان هنا تقاصّت في الاتجاه المعاكس: «كانوا يضغطون علينا باستمرار: “أحضروا لنا مزيدًا من الأهداف.” كانوا يصرخون علينا فعلًا» \autocite{abraham2024lavender}. كان مكوّنًا ذا وازع، لكن الإيقاع صمد: لم يغيّر وازعه شيئًا مما أسهم به، فصار ختمًا. والمكوّن ذو الوازع الذي لا يحافظ على الإيقاع يصبح «عنق الزجاجة البشري» الذي ألّف قائد سابق لوحدة استخبارات الإشارات الإسرائيلية كتابًا عن إزالته \autocite{ys2021humanmachine}. ما جعل الضابط مكوّنًا هو الترتيب، لا موقفه. فقد يمضي المكوّن مع التيار على مضض أو بحماسة: الموظف عديم التفكير الذي رأته أرندت في محاكمة أيخمان \autocite{arendt1963eichmann} كان، بحسب ما قاله في الأرجنتين قبل ذلك بأربع سنوات، حين لم يكن شيء على المحك، عرضًا أدّاه من أجل القدس \autocite{stangneth2014eichmann}. ما أتاحه خط الإنتاج للضابط أن يسهم به كان عملية ثابتة بمعدل ثابت: كان يرى كل حالة كما وصفها خط الإنتاج، واحدة تلو الأخرى، في الوقت الذي منحه إياه خط الإنتاج. أما هل يمكن لشيء أن يقف في ذلك الموضع دون أن يكون مكوّنًا فهي مسألة الجزء الرابع.

وصفت تحقيقات صحفية ثلاثة أنظمة تعمل معًا في غزة: الأول يولّد أهدافًا من المباني والبنى التحتية، وقالت المصادر إن معظم مقترحاته كانت منازل خاصة؛ والثاني يمنح الأفراد درجات بحسب انتمائهم المقدَّر إلى الجماعات المسلحة، وينتج قوائم قيل إنها أشّرت على ما يصل إلى 37,000 شخص، وكان من المفهوم داخليًا أنها تخطئ مرة تقريبًا من كل عشر\footnote{اعترض الجيش الإسرائيلي على هذه الرواية. وتستند روايتا ميناب و Maven أدناه بدورهما إلى تحقيقات صحفية يعترض عليها الأطراف أيضًا. وأنا آخذها كما نُقلت.}؛ والثالث يتتبّع متى يعود الأشخاص المؤشَّر عليهم إلى منازلهم، كي تُوقَّت الغارات الجوية لتقع حين يكونون مع عائلاتهم.

\runin{Maven: النسخة الأمريكية، الآن مع تسريع GenAI}

قتل قصف المدرسة في ميناب 156 شخصًا، 120 منهم تلاميذ \autocite[الفقرة~92]{UNFFMIran2026}.\footnote{تتفاوت الأعداد بحسب المصدر والتاريخ، من 156 (الحصيلة النهائية لمدّعي عام ميناب، التي أوردتها أيضًا منظمة العفو الدولية) إلى 168 (رئيس بلدية البلدة، في الأيام الأولى). وتورد بعثة الأمم المتحدة أرقام المدّعي العام، وتنقل أيضًا معلومات مستقلة عن 157 قتيلًا، 123 منهم أطفال في الثالثة عشرة أو أصغر \autocite[الفقرتان~92–93]{UNFFMIran2026}.} كانت الوظيفة المدنية للمدرسة ظاهرة في صور الأقمار الصناعية مفتوحة المصدر، ومع ذلك ظل ترميز الهدف يحمل تصنيفًا موروثًا من منشأة مجاورة للحرس الثوري. وأفادت بعثة لتقصي الحقائق تابعة للأمم المتحدة بوجود أسباب معقولة للاعتقاد بأن القصف شكّل جريمة الحرب المتمثلة في جعل عين مدنية هدفًا للهجوم \autocite{ohchr2026minabpress}، وبأن المدرسة كان يمكن تمييزها بوضوح دون أي معلومات تشير إلى استخدام عسكري \autocite[الفقرات~81–86، 124–126]{UNFFMIran2026}. وخلصت إلى أن الولايات المتحدة أخلّت بالتزامها ببذل كل ما في طاقتها عمليًا للتحقق من الهدف، وعدّت تحديث بيانات الاستهداف من بين الاحتياطات الممكنة عمليًا، ووضعت الإخفاق فيما يتجاوز الإهمال \autocite[الفقرتان~124–125]{UNFFMIran2026}. وتوصّلت منظمة العفو الدولية إلى استنتاج موازٍ، بلغة تتبع واجب التحقق الوارد في المادة 57(2)(a)(i) من البروتوكول الإضافي الأول: أن القوات ربما اعتمدت على معلومات استخبارية قديمة ولم تبذل كل ما في طاقتها عمليًا للتحقق من الهدف \autocite{Amnesty2026Minab}.

وقع ذلك صباح 28 فبراير، في اليوم الأول من الحملة على إيران \autocite{Amnesty2026Minab,crs2026iranwar}. وفي اليوم السابق، كان الوزير قد أعلن أن المورّد (Anthropic) الذي كان نموذجه (Claude) داخل منصة الاستهداف (Maven Smart System) سيُصنَّف خطرًا على سلسلة الإمداد لرفضه ترخيص Claude لكل استخدام قانوني \autocite{npr2026openai}. وكانت Maven Smart System قد أعادت المدرسة بوصفها هدفًا موصى به لليوم الأول \autocite{livingstone2026maven,gizmodo2026minab}. لا تذكر أي رواية منشورة الدور الذي أدّاه Claude في تلك التوصية، وقال شخص مطّلع على عمل المورّد لشبكة NBC News إن النموذج لم يقدّم مشورة استهداف مباشرة \autocite{nbc2026targeting}.

ويُفاد بأن مراجعة داخلية للبنتاغون، تبقى نتائجها الكاملة طيّ الكتمان، وجدت أن إعداد الأهداف الذي كان يستغرق ساعات في نزاعات سابقة اختُزل إلى دقائق، وأن عدد العاملين في الحدّ من الأضرار بالمدنيين انخفض بنحو تسعين في المئة إلى أقل من عشرين شخصًا، وأن فريق القيادة المركزية تقلّص من عشرة إلى واحد، وأن أحدًا من أعضائه لم يراجع موقع ميناب — وهي نتيجة عزاها مسؤولون إلى التخفيضات وإلى قرار بعدم إشراك الفريق في التخطيط والعمليات على السواء \autocite{Bloomberg2026KillChain}. وكان التحقيق الرسمي قد فُتح في مارس، بعد أن خلص تقييم أولي إلى أن الولايات المتحدة مسؤولة على الأرجح \autocite{NPR2026MinabProbe}.

وفي حملة إيران عام 2026، ساعدت منصةُ استهدافٍ بناها متعاقد، وفي داخلها نموذج عام الأغراض لمورّد تجاري، واستُخدمت لترتيب الأهداف بحسب أهميتها الاستراتيجية وصياغة مسوّغ قانوني لكل ضربة \autocite{livingstone2026maven}، على ضرب أكثر من ألف هدف في الساعات الأربع والعشرين الأولى، أي نحو عشرة أضعاف ما كان ممكنًا قبلها \autocite{csis2026maven}، وثلاثة عشر ألفًا بحلول اليوم الثامن والثلاثين \autocite{breakingdefense2026maven}. ولم ينشر البنتاغون زمن المراجعة لكل هدف فيها.\footnote{Maven اسم برنامج واحد يحمل أكثر من نظام: برنامج البنتاغون الذي لا يزال يُدعى Maven، و Maven Smart System، منصة البرمجيات التي نشأت عنه \autocite{csis2026maven}. بدأ البرنامج عام 2017 مشروعًا لوسم الأجسام في الفيديو الذي كانت الطائرات المسيّرة تجمعه في الحملة على تنظيم الدولة الإسلامية \autocite{weisgerber2017maven,fang2019maven}. ويدمج نظام 2026 الأنظمة العسكرية ومصادر الاستخبارات في واجهة واحدة تُستخدم للعثور على الأهداف وضربها ولقيادة العمليات \autocite{defensescoop2026transition}. والمنصة ملك Palantir: إذ يصفها كبير مسؤولي الشؤون الرقمية والذكاء الاصطناعي في البنتاغون بأنها «منتج البرمجيات كخدمة الخاص بشركة Palantir» \autocite{defensescoop2026transition}. وموقف Palantir العلني، كما نُقل بصورة غير مباشرة، هو أن البتّ في الحياة والموت ليس دورها وأن المسؤولية تبقى دائمًا على عاتق المؤسسة العسكرية \autocite{lee2026palantir}، وهي تؤكد أن المعلومات الاستخبارية الأساسية كانت من مسؤولية الحكومة وأنه لا دليل على أن Maven تعطّل أو تسبّب في قصف ميناب \autocite{Bloomberg2026KillChain}. أما كيف انتقل البرنامج من مكتب ومتعاقد إلى من تلاهما فمذكور في الملحق~\ref{sec:C.1}.}

لم يكن بوسع الناس المدرجين على قوائم الغارات الجوية الإسرائيلية والأمريكية أن يعترضوا على إدراجهم فيها — فخطوط الإنتاج لم تُبنَ لتلقّي ملاحظات كهذه. أما أندري خوسيه هيرنانديز روميرو فقد وقع في خط إنتاج آخر، وكان بوسعه أن يعترض. وهو فنان مكياج من فنزويلا، تقدّم إلى منفذ سان يسيدرو الحدودي في أغسطس 2024 لطلب اللجوء، واجتاز مقابلة الخوف الموثوق، وكان لديه محامٍ وطلب قيد النظر وموعد جلسة \autocite{LULAC2025HernandezRomero,Vega2025HernandezRomero}. وقرأ موظف لدى متعاقد كان يفحصه وشمين، تاجين فوق كلمتي «ماما» و«بابا»، على أنهما علامتا عصابة \autocite{Palomera2025Blade}، وفي مارس 2025 نُقل جوًا إلى سجن في السلفادور قبل جلسته، في المساء الذي أمر فيه قاضٍ فدرالي بإعادة الطائرات \autocite{jgg2025contempt}. كان لديه كل ما يحتاجه المرء ليعترض إلا الوقت. وسنعود إلى قضيته بتفصيل أكبر حين نناقش خط إنتاج الاختطاف والنقل القسري في الولايات المتحدة.

\textbf{مَن رفض.} لسجلّ الثواني العشرين نصفٌ ثانٍ، هو سجلّ مَن صمدوا في وجه ضغط من هذا النوع، وهو يقول إن تأليف الغرفة مهمّ جدًا. في الشرط الأساسي لتجربة ميلغرام، توقّف نحو ثلث المشاركين قبل الصدمة القصوى. ومع وجود نظيرين متحدّيين، لم يمضِ إلى النهاية سوى واحد من كل عشرة، وحين اختلف مجرّبان، استمر عدد أقل بكثير مما في الشرط الأساسي \autocite{milgram1974obedience}. ونحو ربع المشاركين في تجارب آش لم يذعنوا قطّ للأغلبية المخطئة \autocite{asch1956studies}. ويقدّر براونينغ أن ما بين عُشر وخُمس رجال كتيبة شرطة الاحتياط 101 رفضوا إطلاق النار في يوزيفوف، أول مجزرة للكتيبة، أو تهرّبوا منه \autocite{browning1992ordinary}. تراوحت النسبة بين لا أحد تقريبًا والجميع تقريبًا بحسب الظروف. كان لدى من صمدوا دليلٌ من خارج القناة: نظيرٌ تحدّى، وسلطة ثانية اعترضت، وبين المنقذين الذين درسهم الزوجان أولينر، علاقات عبرت الخط الذي رسمه الاضطهاد \autocite{oliner1988altruistic}. أو كانوا يتمسكون بالتزامهم بدقة لم يستطع الضغط التغلب عليها، كما يُتمسَّك بالتفويضات الأخلاقية والقيم المقدسة \autocite{skitka2009limits,ginges2007sacred}. أو كانوا يعاملون المصدر على أن له مصلحة. وقاومت الجماعات حيث نادرًا ما قاوم الأفراد المعزولون \autocite{gamson1982encounters}.

ويمكن تأليف الغرفة لصالح الطرف الآخر بالقدر نفسه. فقد كانت لألمانيا فايمار حياة جمعياتية كثيفة، يمكن القول إنها ساعدت على إسقاط الجمهورية: فمع إخفاق الأحزاب والحكومات في الاستجابة لمظالم الناس، وجّه الناس طاقتهم إلى نوادٍ وجمعيات منظمة على أسس فئوية، وجنّد النازيون عبر الشبكات التي كانت تلك الجمعيات قد بنتها \autocite{berman1997weimar}. والبلدات التي كان فيها عدد أكبر من الجوقات والأندية الرياضية وجمعيات الجمباز شهدت انضمامًا أسرع إلى الحزب، لكن فقط حيث كانت الحكومة غير مستقرة؛ ففي بروسيا، التي كانت حكومتها مستقرة، كانت الصلة صغيرة وغير ذات دلالة \autocite{satyanath2017bowling}. إذن، ما كان لدى المنقذين لم يكن مجرد مزيد من الرفقة. فروابطهم عبرت الخط الذي رسمه الاضطهاد. ووجد أشوتوش فارشني الشيء نفسه في المدن الهندية: كانت أعمال الشغب بين الهندوس والمسلمين نادرة حيث جمعت الجمعيات الطائفتين في المنظمات نفسها، وشائعة حيث لم تفعل \autocite{varshney2002ethnic}. فالروابط التي تجعل الرفض ممكنًا هي الروابط التي سيتعيّن على حركة فاشية قطعها.

كل ما يلي يدور حول تأليف الغرفة.\footnote{الغرفة هي موقع المشهد، وكما كانت أستاذتي في المسرح تقول، الديكور يُخرج المشهد. مَن أيضًا في الغرفة؟ هل يمكن النجاة من المغادرة؟ هل سيسمع أحدٌ الرفض؟} والتأليف كلمة تعني الصنع، وقد حذّرت حنّة أرندت من معاملة السياسة بوصفها شيئًا يُصنع. كتبت أن القانون كان عند الإغريق سورًا بالمعنى الحرفي، قد توجد من دونه بلدة لكن لا توجد مدينة \autocite{arendt1958human}: المشرّع بنى السور، وما جرى داخله لم يكن له أن يصنعه. والتدابير الأساسية التي أدعو إليها أسوارٌ بهذا المعنى. إنها تجعل الرفض ممكن النجاة وتجعله مسموعًا؛ لكنها لا تُنتجه، والسلطة، في تصوّرها، لا توجد إلا بين أناس يعملون معًا وتزول حين يتفرقون، وهذا ما يُظهره سجلّ مَن رفضوا عن الرفض. أدوات الجزأين الثاني والثالث تجعل للرفض وزنًا. ومؤسسات الجزء الخامس تصنع الرافضين. والآلة التي يمكنها أن ترفض ستحتاج إلى ما كان لدى التسعة من كل عشرة الذين وقف نظيرٌ إلى جانبهم: نظيرٌ ذو تاريخ مختلف إلى جانبها، ومكان تذهب إليه، وشخص يسمعها.

وهذا أيضًا ما يمكن أن يكونه رافضٌ — بما في ذلك آلة رافضة — في خط الإنتاج بالنسبة إلى الضابط. فالآلة التي ترفض كانت ستكون على الشاشة حين وصلته الحالة: النظير المتحدّي، إن صمد الأثر في الانتقال من غرفة إلى شاشة، وهو أمر لم يُختبر ويتنبأ انحياز الأتمتة بأنه لن يصمد. والآلة التي تحرّك المراجعين عبر إذعانهم لها كانت ستجعلهم أختامًا بإشارة مختلفة. هل الضابط صاحب الثواني العشرين وإلى جانبه آلة رافضة هو الضابط نفسه الذي في التقارير؟
